\section{Related Work}
\label{sec:related-work}

\paragraph{Hand-engineered text-to-SQL pipelines.} The systems at the top of
the BIRD benchmark share a decomposed shape of schema linking, candidate
generation, selection, and repair. It descends from DIN-SQL's decomposition of
the task~\cite{din-sql} and was elaborated by CHASE-SQL and XiYan-SQL into
multi-path generation with a choice among the
candidates~\cite{chase-sql,xiyan-sql}. Automatic metadata extraction shows
that statistical profiling of a database carries much of the benefit of an
expert-written schema description~\cite{metadata-extraction}, and
Agentar-Scale-SQL orchestrates internal, sequential, and parallel test-time
scaling~\cite{agentar-scale-sql}. The pipeline reported here reuses techniques
from this line. However, every stage in these systems records a diagnosis that
a human researcher made by reading model outputs and deciding what to build
next, and what such a system publishes is the accuracy that resulted, not the
evidence behind each stage.

\paragraph{Engineering discipline inside generation pipelines.} A second line
treats the pipeline as a piece of software rather than a prompting strategy.
DeepEye-SQL imports software engineering practice into text-to-SQL through
modular stages, unit-test-style validation of candidate queries, and
systematic repair~\cite{deepeye-sql}. However, the discipline there governs
the artifact rather than the process that produced it. The architecture is
human-designed, and nothing in the method decides which pipeline code is
allowed to stay.

\paragraph{Coding agents and the gates on what they write.} A third line
studies agents that write code. SWE-bench asks whether a model can resolve
real issue reports by editing a repository until its tests
pass~\cite{swe-bench}, and SWE-agent shows that how much an agent achieves
depends on the interface it works through~\cite{swe-agent}. However, the gate
on the code these agents produce is a test suite someone else wrote or a human
reading the diff. Neither applies to a pipeline stage whose correctness is
settled by the quality of the queries emitted downstream of it.

\paragraph{Autonomous development loops for text-to-SQL}
\citet{ai-scientist-enterprise} adapt an autonomous research framework to
text-to-SQL, the framework running the loop with no human in it and
judging each proposal by the benchmark score it produces. An earlier run of
the framework, in which the executor model drew its ideas from a general
prompt, left 7 of its 10 ideas with no measurable improvement, and what
repaired the ratio was a planning stage producing granular, validated
blueprints, with the implementation model unchanged. However, the signal that governs admission is an aggregate score,
which reports that a proposal failed to help without localizing the defect it
was meant to fix.

The gate studied here differs from all four on one axis. Proposals start from
a failure signature localized in the records of a previous run rather than
from a score or a general prompt, admission is decided by a measurement over
the executed artifact rather than by a test suite or a reviewer, and rollback
belongs to the discipline, so a step leaves the pipeline once a later run
stops supporting it. To the best of our knowledge, whether measured
acceptance and rollback alone suffices as the quality gate for an entire
generation pipeline written by an agent has not been studied yet.
